\section[Cote 113 : l'enveloppe de Karoubi et l'invariance de Morita]{Cote 113 : l'enveloppe de Karoubi et l'invariance de Morita\protect\verifie{Folder113}}\label{sec:113}

La cote 113, \emph{Abelianization : notes manuscrites (s.d.)}, douze pages que l'inventaire date \q{[à partir de 1982]}, construit sous le titre \q{$\otimes$ dans Cat} le produit tensoriel de petites catégories $k$-linéaires, et montre qu'il calcule celui de leurs catégories de modules. La page~3 pose les fondations. Pour $P$ petite catégorie $k$-additive, $P^k$ est \emph{défini} comme $\operatorname{Hom}_{\mathbf Z}(P^\circ, Ab)$, identifié à $\operatorname{Hom}_k(P^\circ, Ab_k)$ avec la parenthèse \q{ne dépend pas du choix de la str.\ $k$-additive}. Le plongement $P \hookrightarrow P^k$ est une \q{inclusion additive pl.\ fid.}, d'où \q{$\operatorname{Kar} P \approx \mathrm{Ul}\operatorname{Proj}(P^{k})$}, les deux lettres devant \q{Proj} n'étant pas identifiées. Une note en marge ajoute \q{$P \to \operatorname{Kar} P$ induit $(\operatorname{Kar} P)^{k} \overset{\sim}{\to} P^{k}$} (\lot{113}{1}{3}). La lecture lit la première équivalence comme $\operatorname{Kar}(P) \simeq \operatorname{Proj}_{\mathrm{tf}}(P^k)$, les projectifs de type fini, et nomme l'ensemble invariance de Morita.

\begin{definition}
Soit $P$ une catégorie. Son \emph{enveloppe de Karoubi} $\operatorname{Kar} P$ a pour objets les couples $(X, e)$ formés d'un objet $X$ de $P$ et d'un idempotent $e = e^2$ de $\operatorname{End}(X)$, et pour flèches $(X, e) \to (Y, e')$ les $f : X \to Y$ tels que $e' f = f = f e$ ; la composition est celle de $P$, et l'identité de $(X, e)$ est $e$. Le foncteur $\iota : P \to \operatorname{Kar} P$ envoie $X$ sur $(X, \mathrm{id}_X)$. Une catégorie est \emph{karoubienne} si tout idempotent s'y scinde. Si $P$ est préadditive, $\operatorname{Kar} P$ l'est aussi, ses groupes de flèches étant des sous-groupes de ceux de $P$.
\end{definition}

\begin{enonce}[Page 3, la note marginale]
Soient $P$ une catégorie préadditive et $\mathcal D$ une catégorie préadditive karoubienne --- par exemple $Ab$, ou la catégorie des $k$-modules. La restriction le long de $\iota^\circ : P^\circ \to (\operatorname{Kar} P)^\circ$ est une équivalence de la catégorie des foncteurs additifs $(\operatorname{Kar} P)^\circ \to \mathcal D$ sur celle des foncteurs additifs $P^\circ \to \mathcal D$. Pour $\mathcal D = Ab$, c'est l'équivalence $(\operatorname{Kar} P)^k \simeq P^k$, avec la définition $P^k = \operatorname{Hom}_{\mathbf Z}(P^\circ, Ab)$ de la page.
\end{enonce}

\begin{lemme}\label{lem:113-univ}
Soient $P$ une catégorie et $\mathcal D$ une catégorie karoubienne. La restriction le long de $\iota : P \to \operatorname{Kar} P$ est une équivalence de la catégorie des foncteurs $\operatorname{Kar} P \to \mathcal D$ sur celle des foncteurs $P \to \mathcal D$.
\end{lemme}

\begin{proof}
C'est la propriété universelle classique de l'enveloppe de Karoubi, complétion libre par les rétractes : un foncteur $G : P \to \mathcal D$ se prolonge en envoyant $(X, e)$ sur l'objet par lequel se scinde l'idempotent $G(e)$, et ce prolongement est unique à isomorphisme unique près. On ne la redémontre pas ici.
\end{proof}

\begin{lemme}\label{lem:113-op}
Il existe un isomorphisme de catégories $\operatorname{Kar}(P^\circ) \to (\operatorname{Kar} P)^\circ$, bijectif sur les objets, qui fait commuter les plongements : le composé $P^\circ \xrightarrow{\iota_{P^\circ}} \operatorname{Kar}(P^\circ) \to (\operatorname{Kar} P)^\circ$ est $\iota_P^\circ$.
\end{lemme}

\begin{proof}
Un objet de $\operatorname{Kar}(P^\circ)$ est un couple $(X, e)$ où $e$ est un idempotent de $\operatorname{End}_{P^\circ}(X) = \operatorname{End}_P(X)$ ; on l'envoie sur l'objet $(X, e)$ de $(\operatorname{Kar} P)^\circ$. Une flèche $(X, e) \to (Y, e')$ de $\operatorname{Kar}(P^\circ)$ est une flèche $f : Y \to X$ de $P$ telle que $f e' = f = e f$, c'est-à-dire exactement une flèche $(Y, e') \to (X, e)$ de $\operatorname{Kar} P$, ou encore une flèche $(X, e) \to (Y, e')$ de $(\operatorname{Kar} P)^\circ$ ; on l'envoie sur elle-même. Ce foncteur respecte identités et compositions, il est bijectif sur les objets et sur chaque ensemble de flèches. Sur $X \in P$, les deux composés donnent $(X, \mathrm{id}_X)$, et sur les flèches l'identité.
\end{proof}

\begin{lemme}\label{lem:113-restr}
Pour $\mathcal D$ karoubienne, la restriction le long de $\iota_P^\circ : P^\circ \to (\operatorname{Kar} P)^\circ$ est une équivalence de la catégorie des foncteurs $(\operatorname{Kar} P)^\circ \to \mathcal D$ sur celle des foncteurs $P^\circ \to \mathcal D$.
\end{lemme}

\begin{proof}
Par le lemme~\ref{lem:113-op}, cette restriction est la composée de la restriction le long de l'isomorphisme $\operatorname{Kar}(P^\circ) \to (\operatorname{Kar} P)^\circ$, qui est une équivalence, et de la restriction le long de $\iota_{P^\circ}$, qui en est une par le lemme~\ref{lem:113-univ} appliqué à $P^\circ$.
\end{proof}

\begin{lemme}\label{lem:113-decomp}
Toute flèche $f : (X, e) \to (Y, e')$ de $\operatorname{Kar} P$ s'écrit $f = p' \circ \iota(f) \circ j$, où $j : (X, e) \to (X, \mathrm{id})$ et $p' : (Y, \mathrm{id}) \to (Y, e')$ sont les flèches données par $e$ et par $e'$.
\end{lemme}

\begin{proof}
Ce sont bien des flèches de $\operatorname{Kar} P$, puisque $\mathrm{id} \cdot e = e = e \cdot e$ et $e' \cdot e' = e' = e' \cdot \mathrm{id}$ ; et le composé est $e' f e = f$.
\end{proof}

\begin{lemme}\label{lem:113-add}
Soient $P$ et $\mathcal D$ préadditives et $F : (\operatorname{Kar} P)^\circ \to \mathcal D$ un foncteur. Si $F \circ \iota^\circ$ est additif, $F$ l'est.
\end{lemme}

\begin{proof}
Par le lemme~\ref{lem:113-decomp}, pour toute flèche $h : (X, e) \to (Y, e')$ de $\operatorname{Kar} P$, on a $F(h) = F(j) \circ (F \circ \iota^\circ)(h) \circ F(p')$, les flèches $j$ et $p'$ ne dépendant que de la source et du but de $h$, et $h$ étant regardée à droite comme une flèche $Y \to X$ de $P^\circ$. Pour deux flèches $h_1, h_2$ de même source et de même but, la somme dans $\operatorname{Kar} P$ est la somme dans $P$ ; l'additivité de $F \circ \iota^\circ$ et la bilinéarité de la composition dans $\mathcal D$ donnent alors $F(h_1 + h_2) = F(h_1) + F(h_2)$.
\end{proof}

\begin{proof}[Démonstration de l'énoncé]
Les foncteurs additifs forment une sous-catégorie pleine de celle de tous les foncteurs. La restriction le long de $\iota^\circ$ étant pleinement fidèle sur tous les foncteurs par le lemme~\ref{lem:113-restr}, elle l'est sur les foncteurs additifs, et elle envoie un foncteur additif sur un foncteur additif. Elle est essentiellement surjective : soit $G : P^\circ \to \mathcal D$ additif ; par le lemme~\ref{lem:113-restr}, il existe $F : (\operatorname{Kar} P)^\circ \to \mathcal D$ et un isomorphisme $F \circ \iota^\circ \simeq G$ ; un foncteur isomorphe à un foncteur additif est additif, donc $F \circ \iota^\circ$ l'est, et $F$ est additif par le lemme~\ref{lem:113-add}. Les catégories $Ab$ et des $k$-modules sont préadditives et karoubiennes, tout idempotent y ayant une image.
\end{proof}

\begin{remarque}
L'énoncé ne demande ni que $P$ soit petite, ni que le but soit abélien. Il suffit que $P$ soit préadditive et $\mathcal D$ préadditive et karoubienne ; l'énoncé vaut en particulier pour les foncteurs additifs à valeurs dans les $k$-modules.
\end{remarque}

\begin{remarque}
L'énoncé est l'invariance classique de la catégorie des modules par passage à l'enveloppe de Karoubi, déduite de la propriété universelle de celle-ci.
\end{remarque}
